\documentclass{article}

\textwidth=6.5in
\textheight=8.5in
\voffset=-54pt
\oddsidemargin=0in
\evensidemargin=0in
\setlength{\parskip}{8pt}
\setlength{\parindent}{0pt}

\usepackage{amsmath, amssymb}
\usepackage{ifthen}

\usepackage[inline]{enumitem}
\setlist{topsep=0pt, parsep=8pt}

\usepackage[rgb,table]{xcolor}\convertcolorsUfalse
\usepackage{graphicx}

% change hyperref options
\PassOptionsToPackage{hyphens}{url}
\usepackage[bookmarks,colorlinks,linkcolor=black,citecolor=blue,pdfstartview=FitH,urlcolor=blue]{hyperref}
\hypersetup{pdfpagemode=UseNone}
\usepackage{bookmark}

\usepackage{graphicx}
\usepackage{multicol}
\usepackage{framed}
\usepackage{array}

\usepackage{icomma}

\usepackage{marginnote}
\usepackage[colorinlistoftodos]{todonotes}
\renewcommand{\marginpar}{\marginnote}

\usepackage{soul}
\usepackage[normalem]{ulem}

\usepackage{pgfplots}
\pgfplotsset{compat=1.18}  
\pgfplotscreateplotcyclelist{mycolor}{black, blue, red, green!70!black, magenta, purple, cyan, orange }  
\pgfplotsset{
  disabledatascaling,
  cycle list=tikzcolor, 
  axis lines=middle,
  every axis plot/.style={no marks, thick},
  every axis/.style={
    enlarge x limits=false,
    enlarge y limits=auto,
    xlabel style={at={(current axis.right of origin)},anchor=west},
    ylabel style={at={(current axis.above origin)},anchor=south},
    % extra x and y ticks for asymptotes
    extra x tick style={grid=major, grid style={dashed,black}},
    extra y tick style={grid=major, grid style={dashed,black}},
    /pgf/number format/1000 sep={},
  },    
  tick label style = {font=\footnotesize, fill=\currentbgcolor, inner sep=0pt},    
  xticklabel shift=1pt,    
  scaled ticks=false,
  scaled y ticks=false,
  unbounded coords=jump,
  samples=101,
  minor tick length=0,
  scatter/classes={
    open={mark=*,draw=black,fill=white, mark size=2, thin},
    closed={mark=*,draw=black,fill=black, mark size=2, thin}
  },
  width=3in,
}
\pgfplotsset{open/.style={only marks, mark=*, fill=white, thin, mark size=2}}
\pgfplotsset{closed/.style={only marks, mark=*, draw=black, fill=black,
    thin, mark size=2}}
\pgfplotsset{labeled/.style={nodes near coords, only marks, mark=*, draw=black, fill=black, thin, mark size=2, point meta=explicit symbolic}}

\usepackage{tikz}
\usetikzlibrary{
  matrix,%
  % enables the snake paths
  decorations.pathmorphing,%
  % enables bigger arrows
  decorations.markings,%
  decorations.pathreplacing,%
  decorations.text,%
  calc,%
  patterns,%
  shapes.geometric,%
  arrows,
  decorations.shapes,
  positioning,
  plotmarks,
  intersections
}
\tikzset{>=latex} % sets default arrow type in tikz
\tikzset{font=\normalsize}
\tikzset{mark size=1.5pt, mark options=thin}
\tikzset{pin distance=4pt,
  every pin edge/.style={<-, thin, shorten <= -2pt}}

\interfootnotelinepenalty=10000
\renewcommand{\thefootnote}{(\arabic{footnote})}

\usepackage{multirow, bigdelim}

% change to \usepackage[sols]{handout} to see solutions
\usepackage[sols, grading, workshop]{handout}

\newcommand\iscurrentenumi[1]{%
  \ifthenelse{\equal{#1}{\theenumi}}{}{#1}}
\setenumerate[1]{ref=\#\arabic*}
\setenumerate[2]{ref=\protect\iscurrentenumi{\theenumi}(\alph*)}
\newcommand\enumiiprefix[2]{%
  \ifthenelse{{\equal{#1}{\theenumi}}\AND{\equal{#2}{(\alph{enumii})}}}{}{}%
  \ifthenelse{{\equal{#1}{\theenumi}}\AND{\NOT{\equal{#2}{(\alph{enumii})}}}}{#2}{}%
  \ifthenelse{\equal{#1}{\theenumi}}{}{#1#2}%
}
\setenumerate[3]{ref=\protect\enumiiprefix{\theenumi}{(\alph{enumii})}\roman*}

\makeatletter

% underline links               
\let\@href=\href
\renewcommand{\href}[2]{\@href{#1}{\ul{#2}}}

\newcommand\calculator{
  \begin{tikzpicture}[x=0.15ex, y=0.15ex, baseline=1.5pt]
    \begin{scope}[thick]
      \draw (0, 0) -- (10, 0) -- (10, 14) -- (0, 14) -- cycle;
      \foreach \x in {10/3, 20/3}
      {
        \draw (\x, 0) -- (\x, 10);
      }
      \foreach \y in {10/3, 20/3, 10}
      {
        \draw (0, \y) -- (10, \y);
      }
    \end{scope}
  \end{tikzpicture}
}
\newcommand{\calcitem}{\refstepcounter{enum\romannumeral\@enumdepth}%
\item[\calculator \csname labelenum\romannumeral\@enumdepth\endcsname]}

\makeatother

\definecolor{purple}{rgb}{0.5,0,1.0}

\newcommand{\doedfinity}[2][\newpage]{
  \begin{problem}[#1]
    Do the Edfinity assignment ``Problem Set #2''.

    We encourage you to write down any scratch work here, as well as
    things you want your future self to remember. Nothing you write
    here will be graded; it's just for you!
  \end{problem}
}

\newcommand{\psrnewpage}{\newpage}
\newcommand{\psreflection}[2][]{

  \ifthenelse{\not\equal{#1}{nocite}}{
    \begin{enumerate}[resume]
    \item
      \begin{problem}[1in]
        Please cite any help you got on this problem set:
        \begin{itemize}[nosep]
        \item If you worked with other students, please list their names here.
        \item If you used resources other than those on our Canvas site, please list them here.
        \end{itemize}
      \end{problem}

    \end{enumerate}
  }

  \probonly{%
    #2

    \psrnewpage

    \emph{Reflection space.} It's a good habit to take a few minutes at the end of each pset to reflect on what you learned and jot down notes to your future self. Here are a few reflection questions to get you started. Nothing you write here will be graded; it's just for you!
    \begin{itemize}
    \item Take a few minutes to look back at the problems. How could you explain to someone else how to do each problem? What was your strategy, and how did you know to use that strategy? If there are problems where you don't feel able to answer this, write down which ones they are. Come back to them in a few days when we post the homework solutions!

      \vspace{1.5in}

    \item Take a look back at the learning objectives on today's worksheet; how are you feeling about each one? If there are ones you know you need more practice with, write those down here.

      \vspace{1.5in}

    \item What did you find most challenging about this material? Do you have any lingering questions about it? If so, write them down here so you don't forget them. At some point, stop by office hours or MQC to ask!

      \vfill
    \end{itemize}      
  }
}

\usepackage{fancyhdr}
\lhead{\textcolor{white}{This content is protected and may not be shared, uploaded, or distributed.}}
\rhead{}
\rfoot{\vspace*{\fill}\footnotesize\textcopyright{} \emph{Harvard Math Ma}}
\renewcommand{\headrulewidth}{0pt}
\pagestyle{fancy}
\begin{document}

\begin{center}
  \large\textsc{Problem Set 18 - More Differentiation Rules, Tangent Line Approximation}
\end{center}

\begin{enumerate}
\item  \note{
    \begin{itemize}[nosep]
    \item Differentiate a function with and without Product Rule (guided).
    \item Differentiate a function with and without Quotient Rule (guided).
    \item Simplify functions using exponent rules and then differentiate (using Constant Multiple and Power Rule)
    \item Exponent review: simplify something like $16^{1/4}$, $16^{3/4}$, and $16^{-3/4}$ (to help with \ref{approximate-cube-root-of-7.5})
    \end{itemize}
  }

  \doedfinity{18}

\item 

  \begin{enumerate}
  \item

    Each of the following functions can be differentiated both with and without the Product Rule. One is easier with the Product Rule, and the other is  easier without. Differentiate them both, using the Product Rule in only \uline{one} part. You need not simplify your answers.

    \begin{grading}
      2 points per part. If they use the Product Rule on both, deduct 1 point. 
    \end{grading}

    \begin{enumerate}

    \item
      \begin{problem}[\vfill]
        $f(x) = \sqrt{x} (3x^5 - \sqrt{7}) \dfrac{1}{x^2}$
      \end{problem}

      \begin{solution}
        First, let's rewrite $f(x)$ as $3x^{7/2} -\sqrt{7} x^{-3/2}$. Now we can
        apply the Power Rule: $f'(x) = 
        3(\tfrac{7}{2}x^{5/2}) -\sqrt{7} (-\frac{3}{2}x^{-5/2})$. 
        
        To simplify, we can factor
        out $\frac{3}{2}x^{-5/2}$ to get 
        $\tfrac{3}{2}x^{-5/2}(7x^5 + \sqrt{7})=\dfrac{3(7x^2+\sqrt{7})}{2x^{5/2}}$.

        If we were to use the Product Rule, we might first rewrite $f(x)$ as 
        $(3x^5-\sqrt{7})(x^{-3/2})$. Then, we'd have
        \begin{align*}
          f'(x) &= (3x^5-\sqrt{7})'(x^{-3/2}) + (3x^5-\sqrt{7})(x^{-3/2})' \\
          &= (15x^4)(x^{-3/2}) + (3x^5-\sqrt{7})(-\tfrac{3}{2}x^{-5/2}).
        \end{align*}
        This looks different from our other answer, but if we factor out 
        $-\frac{3}{2}x^{-5/2}$ and combine like terms, we can see that they are equivalent.
      \end{solution}

    \item
      \begin{problem}[\nstretch{2}]
        $f(x) = (x^5 + 3x^2 -2) (2x^5 - 2\sqrt x + \pi^2)$
      \end{problem}

      \begin{solution}
        We can apply the Product Rule right away.
        \begin{align*}
          f'(x) &= (x^5 + 3x^2 -2)'(2x^5 - 2\sqrt x + \pi^2)
          + (x^5 + 3x^2 -2)(2x^5 - 2\sqrt x + \pi^2)'. \\
          \intertext{Now we can apply the Power Rule.}
          f'(x) &= (5x^4 + 6x)(2x^5 - 2\sqrt x + \pi^2) 
          + (x^5 + 3x^2 -2)(10x^4 - 2(\tfrac{1}{2}x^{-1/2})) \\
          &= (5x^4 + 6x)(2x^5 - 2\sqrt x + \pi^2) 
          + (x^5 + 3x^2 -2)(10x^4 - x^{-1/2}).
        \end{align*}

        If we were to not use the Product Rule, we would first need to expand $f(x)$ as
        \[f(x) = 2x^{10} + 6x^7 -2x^{11/2} + \pi^2 x^5 - 4x^5 -6x^{5/2} + 3\pi^2x^2 
        + 4 x^{1/2} -2\pi^2.\]
        Then we could use the Power Rule.
        \[f'(x) = 20x^{9} + 42x^6 -11x^{9/2} + 5\pi^2 x^4 - 20x^4 -15x^{3/2} + 6\pi^2x 
        + 2 x^{-1/2}.\]
        This looks different from our other answer, but if we expanded our other answer,
        we would see that the two are equivalent.
      \end{solution}

    \end{enumerate}

  \item
    Differentiate the following four functions, aiming for efficiency. \uline{Use the Quotient Rule only twice.} Show your work; you need not simplify your answers.

    \begin{grading}
      2 points each. If they use the Quotient Rule more than twice, deduct 1 point for each additional use.
    \end{grading}

    \begin{enumerate}
    \item
      \begin{problem}[\vfill]
        $f(x) = \dfrac{3x^2 + 7x + \dfrac{1}{x}}{\sqrt[4]{x}}$
      \end{problem}

      \begin{solution}
        We can rewrite $f(x)$ as $3x^{7/4} + 7x^{3/4} + x^{-5/4}$.
        Applying the Power Rule, we get $f'(x) = 3(\tfrac{7}{4}x^{3/4}) 
        + 7(\tfrac{3}{4}x^{-1/4}) + (-\tfrac{5}{4}x^{-9/4}).$
      \end{solution}    

    \item
      \begin{problem}[\vfill]
        $f(x) = \dfrac{\sqrt[4]{x}}{3x^2 + 7x + 1}$
      \end{problem}

      \begin{solution}
        We cannot rewrite the function, so we apply the Quotient Rule:
        \begin{align*}
          f'(x) &= 
          \frac{(\sqrt[4]{x})'(3x^2 + 7x + 1) - (\sqrt[4]{x})(3x^2 + 7x + 1)'}
          {(3x^2 + 7x + 1)^2} \\
          &= \frac{(\tfrac{1}{4}x^{-3/4})(3x^2 + 7x + 1) - (\sqrt[4]{x})(6x+7)}
          {(3x^2 + 7x + 1)^2}.
        \end{align*}
      \end{solution}

    \item
      \begin{problem}[\vfill]
        $f(x) = \displaystyle \frac{(x^2 + 2x)x}{\sqrt{3}}$
      \end{problem}

      \begin{solution}
        We can rewrite $f(x)$ as $\tfrac{1}{\sqrt{3}}(x^3+2x^2)$. Applying the
        Power Rule, we get $f'(x) = \tfrac{1}{\sqrt{3}} (3x^2 + 4x)$.
      \end{solution}

    \item
      \begin{problem}[\vfill\newpage]
        $f(x) = \displaystyle \frac{x}{x^2 + 1} $
      \end{problem}

      \begin{solution}
        We cannot rewrite the function so we apply the Quotient Rule:
        \begin{align*}
          f'(x) &= \frac{(x)'(x^2+1) - (x)(x^2+1)'}{(x^2 + 1)^2} \\
          &= \frac{(1)(x^2+1) - (2x)x}{(x^2 + 1)^2} \\
          &= \frac{-x^2+1}{(x^2+1)^2}
        \end{align*}
      \end{solution}

    \end{enumerate}

  \end{enumerate}

\item \label{quotient-rule}

  In class, we stated the Quotient Rule: $\displaystyle \frac{d}{dx} \left( \frac{f(x)}{g(x)} \right) = \frac{g(x) f'(x) - f(x)g'(x)}{g(x)^2}$.

  However, we didn't see why the Quotient Rule is true. Now, you'll see how to come up with it! Let $Q(x) = \dfrac{f(x)}{g(x)}$. Our goal is to find a formula for $Q'(x)$.

  \begin{grading}
    1 point per part
  \end{grading}

  \begin{enumerate}
  \item \label{quotient-rule-step1}

    \begin{problem}[\vfill]
      We can rewrite the equation $Q(x) = \dfrac{f(x)}{g(x)}$ as $g(x) Q(x) = f(x)$. Differentiate both sides of the equation $g(x) Q(x) = f(x)$. (What rule do you need to use?)
    \end{problem}

    \begin{solution}
      Using the product rule:
      $$\frac{d}{dx} \big( Q(x)g(x) \big) = \frac{d}{dx}
      f(x)$$ $$\fbox{$Q'(x)g(x) + Q(x)g'(x) = f'(x)$}$$
  \end{solution}

  \item

    \begin{problem}[\vfill]
      Remember our goal is a formula for $Q'(x)$. Your answer to
      \ref{quotient-rule-step1} should be an equation involving
      $Q'(x)$, so solve it for $Q'(x)$. 
    \end{problem}

    \begin{solution}
      Our equation becomes: $$Q'(x)g(x) = f'(x) - Q(x)g'(x)$$ $$\fbox{$Q'(x) =\displaystyle \displaystyle\frac{f'(x)-Q(x)g'(x)}{g(x)}$}$$
    \end{solution}

  \item

    \begin{problem}[\nstretch{2}]
      We really want to express $Q'(x)$ in terms of just
      $f(x), f'(x), g(x), g'(x)$. Do this, and simplify your answer until it looks like the Quotient Rule. 
    \end{problem}

    \begin{solution}
      Since $Q(x) = \dfrac{f(x)}{g(x)}$, our previous expression
      becomes
      $Q'(x) = \displaystyle\frac{f'(x) -
        \displaystyle\frac{f(x)}{g(x)}g'(x)}{g(x)}$.

      Let's multiply by $\dfrac{g(x)}{g(x)}$ to deal with the
      fractions in the numerator:
      $$Q'(x) = \displaystyle\frac{g(x)}{g(x)} \cdot
      \displaystyle\frac{f'(x) -
        \displaystyle\frac{f(x)}{g(x)}g'(x)}{g(x)} =
      \fbox{$\displaystyle\frac{g(x)f'(x)-f(x)g'(x)}{g(x)^2}$}$$
    \end{solution}
  \end{enumerate}

\item \label{approximate-cube-root-of-7.5}

  In this problem, you'll approximate $\sqrt[3]{7.5}$.

  Actually, just approximating a quantity often isn't that helpful. For example, if I told you that $\sqrt{250} \approx 15$, you don't have any idea how good of an approximation that is; it might be terrible! It's a little more informative if I tell you something like ``$\sqrt{250}$ is between 15 and 16'' because that gives you some idea of how far off my approximation could be. So, what you'll really do in this problem is to come up with an upper bound and lower bound for $\sqrt[3]{7.5}$.

  \begin{enumerate}
  \item \label{approximate-cube-root-using-tangent-line}

    \begin{grading}
      6 points:
      \begin{itemize}[nosep]
      \item 1 point for choosing to find the tangent line at $x = 8$
      \item 2 points for finding the tangent line correctly: 1 for the slope, 1 for the equation
      \item 1 for correctly plugging in 7.5
      \item 2 points for a correct picture and explanation of why this is an upper bound. 
      \end{itemize}
    \end{grading} 

    \begin{problem}[\vfill\newpage]
      Use an appropriate tangent line approximation to approximate $\sqrt[3] {7.5}$. Express your answer as a simplified fraction. Include a carefully labeled picture to illustrate your approximation, and use your picture to explain whether your approximation is an upper bound or lower bound for $\sqrt[3]{7.5}$.
    \end{problem}

    \begin{solution}  
      We want to approximate the function $f(x) = \sqrt[3] {x}$ near the point $(8, 2)$. To accomplish this task we can use a tangent line approximation. The following is the graph of $f(x)$ along with the tangent line at $x=8$. \begin{center}  
        \begin{tikzpicture}
          \begin{axis}[xmin=-2, xmax=10, cycle list={very thick}, axis
            equal=true] 
            \addplot [domain=0:10, samples=1000]{x/abs(x)*abs(x)^(1/3)}; 
            \addplot+[domain=0:10] [green] {1/12 * (x-8)+2};
          \end{axis}
        \end{tikzpicture}\end{center} 
      The slope of the tangent line at $x=8$ can be calculated by using the derivative. 
      $$ f'(x) = \frac{1}{3} x^{- \frac{ 2}{3} } $$ 
      $$f'(8) = \frac{1}{3} 8^{-\frac{2}{3}} = \frac{1}{3} \cdot \frac{1} {4} = \frac{1}{12} $$ 
      This means that the equation for the tangent line is $L(x) = \frac{1}{12} (x-8) +2$. The big upshot of this is that the function is very well approximated by $L(x)$ for values of $x$ near $8$. 
      $$\sqrt[3]{7.5} \approx L(7.5) =\frac{1}{12} (-.5)+2 = \frac{47}{24} $$
      Note that the tangent line is above the graph of the function so this is an upper bound. 
      $$ \sqrt[3]{7.5} < \frac{47}{24}$$ 
    \end{solution} 

  \item

    \begin{grading}
      3.5 points:
      \begin{itemize}[nosep]
      \item 1 point for knowing to use the secant line between $x = 1$
        and $x = 8$
      \item 1 point for finding the secant line correctly 
      \item 0.5 for correctly plugging in 7.5
      \item 1 points for a correct picture showing this is a lower bound
      \end{itemize}
    \end{grading}

    \begin{problem}[\vfill\newpage]
      If your answer to \ref{approximate-cube-root-using-tangent-line}
      was an \emph{upper} bound for $\sqrt[3]{7.5}$, then use an
      appropriate secant line approximation to find a \emph{lower}
      bound for $\sqrt[3]{7.5}$.

      If your answer to \ref{approximate-cube-root-using-tangent-line}
      was a \emph{lower} bound for $\sqrt[3]{7.5}$, then use an
      appropriate secant line approximation to find an \emph{upper}
      bound for $\sqrt[3]{7.5}$.

      Use a picture to illustrate your work and to show how you know
      that you've found the right kind of bound. 
    \end{problem}

    \begin{solution}  
      Two $x$ values that we know the value of $f(x) = \sqrt[3]{x}$ for are $x=1$ and $x=8$. Below is a graph of the function with the secant line through the points $(8, 2)$ and $(1, 1)$ is drawn. \begin{center} 
        \begin{tikzpicture}
          \begin{axis}[xmin=-2, xmax=10, cycle list={very thick}, axis
            equal=true] 
            \addplot[domain=0:10, samples=1000]{x/abs(x)*abs(x)^(1/3)}; 
            \addplot+[domain=-0:10] [green] {1/7 * (x-8)+2};
          \end{axis}
        \end{tikzpicture}\end{center} 
      The slope of the secant line is $\frac{1}{7}$. This means that the equation for the secant line is $S(x) = \frac{1}{7}(x-8) +2$. 
      $$\sqrt[3]{7.5} \approx T(7.5) = \frac{1}{7} (7.5 -8)+2=2 - \frac{1}{14} = \frac{27}{14}$$
      On the interval $[1, 8]$ the secant line is below the graph of the function, so this is a lower bound, \[\sqrt[3]{7.5}>\frac{27}{14}.\] 

    \end{solution} 

  \end{enumerate}

\calcitem % data from https://www.census.gov/quickfacts/fact/table/MA/PST045222 and https://www.census.gov/quickfacts/fact/table/MA/INC110221

  \begin{grading}
    5 points:
    \begin{itemize}[nosep]
    \item 1 for recognizing that they should use Product Rule
    \item 2 for translating the given information appropriately (like 85000 is the value of a function and 1350 is the value of its derivative). If they get the sign of 25000 wrong, deduct 0.5. 
    \item 1 for units of their final answer
    \item 1 for having an explanation that you can follow
    \end{itemize}
  \end{grading}

  \begin{problem}[\vfill]
    According to the U.S. Census Bureau, in 2022, the average annual income per person in Massachusetts was about \$85,000 and was increasing at a rate of about \$1350 per year. Also, in 2022, the population of Massachusetts was about 7 million people and was decreasing by about 25,000 people per year.

    Gross state income is the total amount of money earned per year by all of the people in the state. Use an appropriate differentiation rule to estimate the rate at which the the gross state income of Massachusetts was changing (with respect to time) in 2022. Please include units with your final answer, and be sure to say whether the gross state income is increasing or decreasing. Explain your reasoning clearly, and say which differentiation rule you're using. 
  \end{problem}

  \begin{solution}
    Let the average annual income per person in Massachusetts in year $t$ be $A(t)$, and
    let the popoulation of Massachusetts in year $t$ be $P(t)$. Then the gross state
    income in year $t$ would be $G(t)=A(t)P(t)$. Using the Product Rule gives us
    $G'(t) = A'(t)P(t) + A(t)P'(t)$. We can then calculate $G(2022$, with units:
    \begin{align*}
        G(2022) &= A'(2022)P(2022) + A(2022)P'(2022) \\
        &= (1.35\times 10^3\text{ dollars per person per year})(7\times10^6\text{ people}) \\
          &\qquad + (-2.5\times10^4\text{ people per year})(8.5\times10^4
          \text{ dollars per person}) \\
          &= (9.45\times10^9\text{ dollars per year}) + (-2.125 \times 10^9
          \text{ dollars per year}) \\
          &= 7.325 \times 10^9\text{ dollars per year}.
    \end{align*}
    Since $G(2022)$ is positive, this means that the gross state income was increasing
    in 2022.
  \end{solution}

\end{enumerate}
\psreflection{For next class, read \href{https://activecalculus.org/single/sec-3-1-tests.html}{Boelkins \S 3.1}.}
\end{document}
